% ================================================================
%  第 11 章 孤独的历险 (原书 p.219-230, 边码 p.122-p.139)
%  源文: _work/src1/task-ch11.md
% ================================================================
% ---------------- p.219 ----------------
\chapter{孤独的历险}

\section{(46) 禁果}

\origpage{122} 记笔记的事,我不得不停了两天。仔细重读之后,我看先前那套戏路子大致上确实是实情的写照——只是这写照现在还得再往深里挖一挖。我尤其该把冥想与数学这两匹“马”各自的长处圈得更准些;也该试着弄明白:究竟是哪些事件、哪些机缘,最终扣动了老板押注的“翻转”——明明有惰性之力拽着他,要他没完没了地守着一注,哪怕那注已经输了。

也许还该探一探孩子自己的偏爱。如今这点已经说定:他时不时就想换一换局,而老板显然也留着最起码的一点伸缩,不至于硬逼着他总玩这个、从不玩那个。这几年,老板学会了把孩子算进去,学着跟他将就,不必等到锅子炸开。这还不是完全的和谐,但已不再是战争,更像是某种诚挚的协约——偶发的紧张往往让它松动,而不是变硬。

只要不被硬顶得太狠,孩子的偏爱天性上是相当柔软的。(这跟老板可不一样——他到头来学到的那一丁点柔软,全是违着自己的意愿、到了晚年才学会的……)不过孩子柔软,并不等于他没有偏爱,不等于他不会对这一样东西的兴致比对那一样更浓。

要把这里头看清,常常一点也不容易:孩子的欲望、老板的偏爱,乃至老板一劳永逸定下的东西,三者之间很难分清。从前我对自己说:冥想比数学更好、更重要、更严肃,如此等等,理由如此如此(不用问,都是最站得住的理由)——那是老板事后给自己找的好理由,好让自己相信他下的那一注确实“押对了”。孩子从不说这样东西比那样“更好”“更重要”。他不爱讲话。他想做什么,只要没人拦,他就去了,从不问这事“重要”还是“更好”。他的兴致,对这件事对那件事、这一刻下一刻,有强有弱。要察觉他的偏爱,听老板那些解释性的话毫无用处——他自称代表孩子说话,其实只能代表他自己。只有观察孩子在玩什么,或许才勘得破他的偏爱。而且就算这样也不那么容易:他兴冲冲地玩着这一样,并不总意味着,要是老板不伸他那根手指头去拨一拨,他就不会欢天喜地地去玩另一样。

% ---------------- p.220 ----------------
\origpage{123} 显然,比其他一切更吸引他的,是未知——是追进夜里那些含混不清的褶子中去,把未知的东西带到光天化日之下,他自己的未知,也是所有人的未知。我感觉到,我说“也是所有人的”的时候,这确实是孩子的欲望,而不是老板的虚荣——那种想叫满堂看客、也叫自己吃一惊的虚荣。还有一点也是说定了的:孩子每回从阁楼、从取之不竭的地窖的幽暗中带回来的,都是些“显然”的、孩子气的东西。它们越是显得显然,他就越高兴。要是它们不显然,那就是他没把活干到底,在幽暗与天光之间半路上停下了。

\origpage{124} 在数学里,“显然”的东西,也是那种迟早总有人会撞上的东西。它们不是“发明”,可以有人做、也可以没人做。它们本来就在那里,从来就在;人人都在它们边上走来走去而不留意,顶多绕着它们兜个大圈子,或者每次都绊一跤、从它们身上跨过去。过了一年,或者过了一千年,难免有人终于留意到那东西,围着它挖,把它刨出来,翻来覆去地看,擦干净,末了给它起个名字。这类工作——我偏爱的工作——每一回换个人也能做;而且,换个人早晚也非做不可\fnoteP{1}{(44)}{不用说,我在这里把一个假设搁在一边了——至少可以说,这个假设绝非不可能:一场原子战争,或诸如此类的另一种“乐子”,猝然闯了进来,一举而永久地终结了那场叫做“数学”的集体游戏,顺带也终结了许多别的东西……} (44)。

在我自己的发现这件事上就完全不一样了——在这场绝非集体的游戏“冥想”里。我所发现的东西,世上没有任何别的人,今天也好、任何别的时刻也好,能替我发现。它只属于我一个人去发现,也就是说,也属于我一个人去承担。这个未知,并不因为事理之必然、几乎不管我肯不肯费心去理会它,就注定要被知道。它在沉默里等着被知道的那一刻;有时时候到了,我听见它在呼唤——而蒙召去认识它的,只有我一个人,只有我里面的那个孩子。这不是一个“缓期执行”的未知。当然,我尽可以听从它的呼唤,也可以躲开,说“明天”,或者说“总有一天”。但那呼唤是冲着我、也只冲着我来的;除我之外没人听得见,除我之外也没人跟得上。

每当我听从了这呼唤,“企业”里就有什么东西变了,多多少少。效果是即时的,当下就被感到是一种恩惠——有时是一种突如其来的解放,一种巨大的宽释:卸下了一副我常常连自己都没察觉在扛着的重担,而它的真实,恰恰由这宽释、这解放显出来。在振幅小一些的音阶上,这类体验在一切发现的工作里都很常见,我也曾有机会谈到过。然而有一件事,使自我发现的工作(不论它做在光天化日之下,还是留在地下)区别于任何别的发现工作:它真的改变了“企业”本身里的什么东西。这不是量上的改变,不是产量的增加,不是产品在尺寸上、甚至在品质上的差别。这是老板与工人—孩子之间关系的改变。也许连老板自己也变了——如果这话除了就他与工人、与孩子的关系而言,还能有别的意思。比如他也许不那么盯着产量了——但这也仍是他与工人关系的一面,因为一种关切、也许一种敬意出现了,而那从前对他是陌生的。凡我冥想过的场合,变化都朝着老板与工人之间关系变得清朗、变得平静的方向。除了某些冥想停留在表面的情形——那是在某种当下而有限的需要压迫之下做的“应景”冥想——那清朗一直延续到今天,平静也是。

% ---------------- p.221 ----------------
\origpage{125} 这就使自我发现的工作带上了别的发现工作所没有的一层意义,尽管许多本质的方面是共通的。在自我的认识里,在自我发现的工作里,有一个维度,使它们区别于一切别的认识、一切别的工作。也许这就是知识之树上那颗“禁果”。也许冥想在我身上施展的那种魔力——更确切地说,是冥想向我揭示其存在的那些奥秘的魔力——正是禁果的魔力。我跨过了一道门槛,恐惧在那里消失了。认识的唯一障碍是一种惰性,一种有时大得可观、但终归有限的惰性,绝非不可逾越。这惰性,我几乎每走一步都感到它,阴魂不散,无处不在。它有时把我惹得冒火,却从没让我泄气。(数学工作里也一样——那里主要障碍也是它,只是分量轻得没法比。)这惰性成了游戏的基本配料之一;更确切地说,成了这游戏里的一位主角——这游戏精巧,而且全不对称,其中有两位——更确切地说三位主角:一边是纵身一跃的孩子,和那(化作惰性的)老板,他竭尽全力地刹车(一面还声称自己不在场);另一边是那美丽未知者隐约可见的身影,满是神秘,既近又远,一面躲闪,一面召唤……

\section{(47) 孤独的历险}

“冥想”在我身上曾经有极大的魔力——跟从前“女人”的吸引力一样大,而它似乎已经取代了后者。我刚写下“曾经”,这并不表示这魔力今天已经熄灭。这一年我把自己投进数学里,它只是退到了背景中。经验告诉我,这局面一天之内就可能翻过来,正如这局面本身也正是一次全然未料的翻转的结果。其实,在我经历过的四个漫长的冥想时期里(其中一个延续了将近一年半),每一次我都理所当然地以为,我会顺着这股势头一直走到最后一口气,尽可能深入地去探生命之谜、人之存在之谜。当笔记堆成吓人的一摞摞、眼看要把我的工作室淹掉时,我甚至终于请人量身定做了一个柜子来安放它们,而且留得很宽(凭一个快速的算术级数推算),好把往后几年里很快会添进去的也一并安放;要是我没记错,我留了大约十五年的余量(这已经相当可观了!)。这一回老板事情办得漂亮,论管家,这是漂亮的管家功夫!这一手,加上把所有与冥想工作或近或远沾边的私人文件来一次大刀阔斧的整理,是他最后着手、并且(几乎)圆满完成的一件差事——就在偏爱与押注翻转之前。这不禁让人琢磨:他是不是心里早有另一层打算,是不是已经看见一册册《数学反思》塞满了那些空着的格子——那格子据说本是留给日后的“笔记”的。

当然,冥想之爱、自我发现之爱,辽阔得足以填满我这辈子直到最后一天。数学之爱也确实没有被耗尽,但也许这份饥渴在往后几年里终会吃饱。我里面有什么在这样盼着,并且觉得数学是一道羁绊,拦着我走那场只有我一人能走的孤独历险。而且我觉得,我里面这个“什么”不是老板,也不是老板的哪一桩心血来潮(按本性,他是分裂的)。我觉得数学之爱还带着老板的印记;无论如何,跟着它走,只是让我的生命在一个封闭的圈子里打转;一个轻省的圈子,一种惰性的运动,绝不是更新。

% ---------------- p.222 ----------------
\origpage{126} 我曾追问:数学之爱在我生命里这样执拗地持续,究竟是什么意思。跟着它走的时候,它并不真的填满我的生命。它给人快乐,也给人满足,但它本身并不足以给人真正的舒展、一种饱满。像一切纯粹智力的活动一样,强烈而长期的数学活动,效果毋宁说是使人变钝。我在别人身上看到这一点,尤其在我自己身上——每当我重新扑进去的时候。这活动太零碎了,它只动用我们直觉、感受诸能力中极小极小的一部分,那些能力由于久不使用而变钝。长期以来我并不觉察,而显然我的大多数同行,当年也不比我更觉察。我想,正是从我单单做冥想以来,我才留意到这件事。稍一留意,它刺眼得很——大剂量的数学使人愚钝。即便在两年半前的那场冥想之后——那场冥想承认了数学之爱确实是一种爱,是我生命里重要的一件事——如今我献身于这爱时,仍留着一份保留,一份迟疑,不是完全的献出。我知道,所谓“完全的献出”其实会是一种退位,会是顺着一种惰性走,会是一种逃避,而不是献出。

对冥想,我身上没有丝毫这样的保留。我投身其中时,是全身心投入,这献出里没有一点分裂的痕迹。我知道,在献出的时候,我与我自身、与世界完全一致——我忠于我的本性,“我就是道”。这份献出对己对人都是好的。它把我向自己敞开,也向他人敞开,用爱去解开我里面还结着的结。

\origpage{127} 冥想把我向他人敞开,它有力量解开我与他的关系,哪怕对方仍旧结着。可是,要找到机会与人多少谈一谈冥想的工作、谈一谈这工作让我认识的这样那样的事,是极难得的。这绝不是因为那些事“太私人”。打个不贴切的比方:对某个时候让我着迷的数学,我只能跟一位具备必备行头的数学家谈,而且他还得同时乐意对之发生兴趣。有时一连几年,我被某些数学东西迷住,却碰不上(甚至不去找)另一个能与之谈论它们的数学家。但我很清楚:只要我去找,就会找到;即便找不到,那也只是运气或时机的问题;让我感兴趣的东西,不可能不使某个人、甚至某些人感兴趣,是十年后还是一百年后,说到底无关紧要。正是这一点,赋予我的工作以意义,哪怕这工作是在孤独中做的。假如世上再没有别的数学家,而且将来也不会再有,我不相信做数学对我还会有意义——我猜,对任何别的数学家、任何别的“研究者”,情形也一样。这就接上了前面那个观察:对我来说,“数学的未知”是还没有任何人知道的东西——它不取决于我一个人,而取决于一个集体的实在。数学是一场集体的历险,已经延续了几千年。

% ---------------- p.223 ----------------
\origpage{128} 至于冥想,要谈论它,“行头”的问题根本不存在;至少到我这个地步不存在,我怀疑它将来也不会存在。唯一的问题,是别人身上有没有一份兴趣,能与我身上这份兴趣相应。所以那是一种好奇心,对着自己和别人身上真正在发生的事——越过那些必不可少的门面;只要人真的想看看门面底下盖着什么,那些门面就藏不住多少东西。但我已经知道,一个人身上出现这种兴趣的时刻,“真相的时刻”,是稀少而短暂的。当然,常能碰到所谓“对心理学感兴趣”的人,读过弗洛伊德、荣格,还读过许多别人,巴不得来几场“有意思的讨论”。他们随身带着那份行头,或轻或重,就是人们所谓的“文化”。它是他们对自己那幅形象的一部分,又加强那幅形象,而那幅形象他们是从来不肯去查看的——跟另一位对数学、对飞碟或者对钓鱼感兴趣的人一模一样。我刚才想说的,不是这种“行头”,也不是这种“兴趣”——尽管同样的字眼在这里指的是性质不同的东西。

\origpage{129} 换一种说法:冥想是一场孤独的历险。它的本性就是孤独的。不仅冥想的工作是一种孤独的工作——我想,一切发现的工作都是如此,哪怕它嵌在一项集体工作里。但由冥想的工作诞生的认识,是一种“孤独”的认识,一种无法分享、更无法“传达”的认识;或者说,它若能被分享,也只在稀少的瞬间。这是一种工作、一种认识,它与最根深蒂固的共识逆向而行,它让每一个人、让所有人不安。这认识当然表达得简单,用些简单清亮的词。当我向自己表达它时,我在表达中学习,因为表达本身就是工作的一部分,由一股强烈的兴趣托着。可是这些同样简单清亮的词,一旦撞上冷漠或恐惧的紧闭的门,就无力向别人传达任何意思。就连梦的语言——气力完全不同,资源无穷,由一位不知疲倦、心怀善意的“做梦者”不断更新——也穿不过那些门……

没有哪一种冥想不是孤独的。只要有一丝想得到谁认可的操心,想得到证实、得到鼓励,就没有冥想的工作,也没有自我的发现。有人会说,任何真正的发现工作,在工作进行的那一刻,也是如此。不错。但在工作本身之外,他人的认可——无论来自亲近的人、同行,还是自己所属的整个圈子——这认可对于这工作在一个献身者生命中的意义,是重要的。这种认可、这种鼓励,是最有力的激励之一,正是它们使“老板”(还是用这个形象)毫无保留地开绿灯,让孩子尽情尽兴。主要正是它们决定了老板的投入。我自己在数学上的投入,情形也一样——那是被嘉当、施瓦茨、迪厄多内、戈德芒这些人,以及他们之后的另一些人的善意、温暖与信任鼓励起来的。相反,对冥想的工作,丝毫没有这样的激励。这是工人—孩子的一种爱,老板心底里还算厚道,多多少少容着它,因为它“不带来”任何东西。它当然结果,但那不是老板巴望的果子。只要他不在这件事上自己骗自己,事情就很清楚:他不会往冥想里投钱——老板天性合群!

只有孩子,天性上是孤独的。

\section{(48) 馈赠与接纳}

昨天谈到冥想的孤独本性时,有个念头掠过:我写了将近六个星期的这些笔记,末了竟成了某种冥想,可它们却是要拿去发表的。这当然必定以许多方式影响了这冥想的形式,尤其是那份求简洁的用心,还有那份求谨慎的用心。冥想的一个本质方面,即对工作当下我身上发生之事保持持续的留意,只偶尔显现,而且浮在表面。这一切肯定影响了工作的进程和品质。但我仍感到它有冥想的品质,首先凭它果实的性质,凭一种对我自己的认识浮现了出来(在此情形下,尤其是对某一段过去的认识),而那是我此前一直回避的。另一个方面是自发性:近五十个“节”或“段”,反思是自发地聚成这些的,而每一个,我在开头时都说不出它的实质会是什么;每一次,实质都只在途中才显露;每一次,工作都带出一些新的事实,或者以新的光照明了此前被忽略的事实。

% ---------------- p.224 ----------------
\origpage{130} 这工作最直接的意义,是与我自己的对话,因而是一场冥想。然而,这场冥想是要发表的,而且还要充当随后那部《数学反思》的“开场”——这一点绝不是可有可无的处境,绝不是在写作过程中不起作用的东西。对我而言,它是这工作意义的一个本质部分。昨天我暗示老板肯定能在其中“捞到他的好处”(他是“捞好处”的老手,什么都少不了他那一份!),这丝毫不意味着它的意义就缩到那一点上——缩成那匹著名的三脚马迟来的、几乎算死后才发生的“回归”!也不止一次,我感到一个行为深层的意思,有时会超出促成它的那些动机(显的或隐的)。而在这“回归数学”里,我猜度着还有另一层意思,不只是某些心理力量相加的结果——那些力量在某个时刻、出于某些缘由,在我这个人身上碰到了一起。

这场我正在进行的“冥想”,是要献给我在数学世界里认识并爱过的人——如果我感到它是那份隐约见到的意义的重要部分,那并不是因为我在盼着这馈赠被人接纳。它被不被接纳,不取决于我,只取决于它所寄予的那人。它被接纳,我当然绝非无所谓。但那不是我的责任。我唯一的责任,是在我所作的馈赠里是真诚的——也就是说,也是做我自己。

冥想让我认识的,是些卑微而显然的事,一些不登大雅之堂的事。也是那些我在任何书里、任何论著里都找不到的事——不论那书多么博学、多么深刻、多么天才——那些没有人能替我找到的事。我盘问过一团“雾”,我费心去听它,我学到了一点卑微的真相,关于一种“竞技态度”及其显然的意义——在我与数学的关系里,也在我与他人的关系里。就算我把《圣经》、古兰经、奥义书都“照原文”读过,再加上柏拉图、尼采、弗洛伊德和荣格,我是个学问广博精深的天才——这一切也只会让我离那真相更远,那是孩子气的、显然的真相。就算我把基督的话重复一百遍——“像小孩子那样的人有福了,因为天国是他们的”——再细细给它作注,那也不过仍旧使我离我里面的孩子更远,离那些让我不舒服、只有孩子才看得见的卑微真相更远。就是这些东西,是我所能献出的最好的。

而我很清楚:这样的话被说出来、被献出去,用的是简单清亮的词,却并不因此就被人接纳。接纳,并不只是收到一条信息,不管是尴尬地收到,还是甚至感兴趣地收到:“嘿,这谁想得到……!”或者:“说到底也没那么稀奇……”。接纳,常常是在献出者身上认出自己。是借着别人的这个人,与自己相识。

\section{(49) 一条裂缝的确认}

这段关于当前工作意义、关于馈赠与接纳的短小思考,像是反思脉络里的一段离题;或者更像是例证,说明某些方面——它们使“冥想”区别于一切别的发现工作,尤其是数学工作。昨天我明白了:这些方面有双重效果,也就是两种方向相反的效果:对“孩子”是一种独一无二的魔力,对“老板”则是全然的无趣。看来这双重效果就在事物的本性里,绝不可能靠什么妥协、什么安排来缓和。不管怎么做,孩子一跟从他真正的偏爱,老板就在里头捞不到他的好处——一点也捞不到!

% ---------------- p.225 ----------------
\origpage{131} 毫无疑问,已经发生的那次翻转,意义就在这里;它很可能在往后几年里把冥想从我生命中一扫而空(除了像三个月前那样的“应景冥想”)。我并不认为,因此那几年就会完全荒芜,正如过去这一年也并不荒芜。但同样确实的是:我在这一年里学到的东西(数学以外),跟此前四年中任何一年学到的东西相比,都微不足道。怪的是,我经历过的四个漫长冥想时期,每一个都是极其饱满的时光,没有任何迹象让人怀疑我心里还有什么落空的东西。然而,锅子既然炸了,就说明某处有压力,而这压力不会是当天才有的;它一定已经存在,在我看不见的某处,存在了几个星期或几个月,而那时我正专注而全然地沉在冥想里。

\origpage{132} 不过这里我是被笔头的冲劲带着走了(更确切地说,是打字机的冲劲)。实情是:(除了最后一个冥想时期——它正冲到半途,就被一连串事件与情势截断了)冥想的强度从某一刻起渐渐低落,就像一道浪,而后面另一道浪正等着接它的位置……说实话,饱满之感也随着这同一运动,只是有个分别:它只在“冥想浪”的时候在场,不在“数学浪”的时候在场。

我试图圈定的这个局面,在我看来已经不再是冲突的局面,但看来它里面仍旧含着冲突的胚芽、冲突的可能。如今它对我来说,也许是凭它对我生命进程的冲击,我内部分裂的最显眼的迹象。这分裂不是别的,就是老板与孩子的分裂。

我无法把它了结。既然它在这一处表现里已经被勘破,我能做的,只是留意它,在往后的月份和年份里追索它的迹象与演变。也许这份对数学的爱——不得不说,有点不合时宜——会在燃烧中烧尽(像我身上另一份爱已经烧尽那样……),让位给那唯一的爱:对自我的发现,对我命运的发现。

这份爱足够辽阔,我说过,足以填满我的生命——而我的整个生命肯定也不够把它耗尽。

\section{(50) 一段过去的重量}

我把《收获与播种》最后收尾完工,已经是几天前的事了——此前一个多月,我一直以为再过几天就要完工。这一回也是,在“收尾”之后,我仍不完全确定自己是否真的完工了——因为有一个问题我确实悬着没解决。那就是“弄明白:究竟是哪些事件、哪些机缘,最终扣动了‘老板’押注的‘翻转’”,使数学取代了冥想的位子,而且是与巨大的惰性之力相抗。这几天,我的思绪不请自来地、相当执拗地回到这个问题上——而这几天我其实已经开始接上另一些全然不同的事,包括一些数学问题(共形几何方面的)。不如再借一借这股冥想的“末势”,多少往深里挖一挖,把地方腾干净。

% ---------------- p.226 ----------------
当我试着“凭感觉”回答为什么“我又弄起数学来了”(指的是那种大规模的、打算长期干的投入,至少以几年计),好几个联想就冒出来。其中最有力的,也许要数我近六七年来在教学活动中终于感到的那种慢性落空之感。就是那种越来越强的感觉:我“大材小用”;而且常常是,我把力气、把我最好的东西投给一些闷闷不乐的学生,而他们对我能给的东西毫不在意。

\origpage{133} 我到处看见值得做的妙事,而它们只等着人去做。往往,一点可怜的行头就足够上手;倒是这些东西自己告诉我们:该发展出什么样的语言来圈住它们,该置办什么样的工具去挖它们。我没法不看见它们——只要由于教学活动而与数学保持着经常的接触(不论那层次多么低),哪怕在我对数学兴趣最边缘的那些时期,也是这样。每一件瞥见的东西背后,只要稍一发掘,又有别的好东西,遮住一些、又露出一些……不论是在数学里还是在别处,只要带着真正的兴趣把目光放上去,就看见一种富饶显露出来,一种深度展开,叫人猜出它们取之不竭。我说的那份落空,就是怎么也没法把这种富饶之感——深之感——哪怕一星半点传给学生:哪怕只是一点想至少把他们手边的东西摸个遍的兴致,在他们无论如何都打算投进所谓“研究”活动的那几个月或几年里,尽情尽兴地干一场——反正也是为了拿这个那个文凭。十年来我教过的学生里,除了两三个,看起来“尽情尽兴”这个念头本身就让他们害怕;他们宁愿几个月、几年地垂着手原地蹭,或者苦哈哈地干一份鼹鼠式的活,那活的来龙去脉他们一概不知,只要末尾有个文凭就行。关于这种创造力的瘫痪,有很多可说的,它与“天赋”或“才能”的有无毫无关系——这又接上了我反思的最开头,那里我曾顺带触及这类阻塞的深层原因。但这不是我在这里要说的话;我这里要说的,毋宁是确认一种慢性的落空状态——这些情形,在最近七年教学活动中不断重演,终于在我身上造成了它。

\origpage{134} “解决”这样一份落空,最现成的办法——至少就它是我身上那位“数学家”的落空、而不是那位教师的落空而言——就是我自己去动手做那些事,至少做其中一部分;那些事我曾绝望地盼着哪个学生终于有一天会接过去做。这也正是我多少在这里那里做过的事:或者借教学活动之余、借教学活动的机会,花几个小时、甚至几天作一次偶然的思考;或者在数学上大发作的时期(那有时来得像真正的爆炸……),一发作就是几个星期、几个月。这样偶发的、一阵一阵的工作,往往只能对一个问题作最初的一遍粗加工,只能给出最零碎的视野——不如说,那是对尚待进行的工作看得更清楚一些;而那工作本身始终搁在那里等着做,而且正因为看得更清,就显得更烫手。两个月前,我给出了一份总体草图,涉及那些我才刚刚开始约略摸到门径的主要题目。那就是《纲领概要》(Esquisse d'un Programme)——我此前已有机会提到过它,它最终将附在这份反思之后,二者合起来构成《数学反思》的第一卷。

相当清楚:单靠这份(不妨说“私下的”)勘察工作,不足以化解我的落空。“大材小用”这个感觉,无疑翻译出一种想有所行动的欲望(我想,它源于自我,也就是说,“老板”的欲望)。这里说的行动,与其说是对他人的行动(比如对我的学生:推动他们,向他们“传达点什么”,或者帮他们拿到某个文凭,好去申请某个职位,等等),不如说是“数学家的”行动:为发现某些无人料到的事实、为某个理论的绽出出一份力,等等……这立刻接上了前面那个观察:数学是一场“集体的历险”。

% ---------------- p.227 ----------------
\origpage{135} 如果我追问自己近十年来做数学时的心境——那是我生命中一段根本不会想到有朝一日还能重新发表东西的时期,而且当时也多少清楚:我眼下或未来的学生,没有一个会用到我的勘察工作——我立刻就看到:那绝非某个只为个人乐趣做事、或者被一种只关乎自己而与他人无关的内在需要所驱赶的人的心境。我做数学的时候,我相信我心里某处是明白的:这些数学是做来传达给别人的,是要成为一件更广大之物的一部分的,我为它出力,而那东西绝不是个体性的。这“东西”,我可以称之为“数学”,或者更好的说法,“我们对数学事物的认识”。这里“我们”一词,首先无疑具体地指我所认识、与我共有某些关切的那些数学家构成的团体;但毫无疑问,它也超出这个狭小的团体,正如它超出我个人。这个“我们”指的是我们的族类——就这一族类通过它在各个时代的某些成员,曾经并且正在关切数学对象世界的诸般实在而言。直到此刻我写下这几行之前,我从未想到过我生命中存在着这“东西”,更没想到要追问它的本性,追问它在我作为数学家与教师的生活里起什么作用。

我提到的那个想有所行动的欲望,在我身上、在我作为数学家的生活里,似乎采取这样的形式:把所有人都不认识的东西从暗处取出来——不只是我不认识的东西(如前所见)——而且,还要把它交给所有人支配,也就是去丰富一份共同的“家产”。换句话说,这是一种欲望:为那超出我个人的“东西”、那份“家产”的扩大与丰厚出一份力。

这份欲望里,当然不缺借自己的作品来扩大自己的欲望。从这一面看,我又碰上了那种对“成长”、对扩大的饥饿——那是自我、是“老板”的特征之一;那是他侵占性的一面,推到极处便是毁灭性的一面(参见注44,第13.1.1节,原书p.260)。不过我同时也明白:想增加那些(或久或暂)多少会挂上我名字的东西的数目,这欲望远不足以穷尽、远不足以覆盖那股更广大的欲望或力量——它推着我去为扩大一份共同家产出力。我觉得这样一份欲望本来可以得到满足(如果不是“在我的企业里”——那里老板还相当有侵占性——至少在一个更成熟的数学家身上),而他自己这个人的角色仍是无名的。这也许就是自我扩大倾向的一种“升华”形式:通过与一件超出自身之物的认同。除非这类力量本身并非源于自我,而是属于一种更细腻、更深的本性;除非它表达的是一种深刻的需要,独立于一切制约,见证着一个人生命与整个族类生命之间那深深的联结——那联结是我们个体存在意义的一部分。我不知道,而追究这样一些如此宽泛的问题,也不是我在这里要说的话。

我要说的,毋宁是(以一种更朴素的角度)考察一个与我这个人有关的具体局面:一个落空的局面,靠零星断续的数学活动得到局部的、暂时的排遣。这样,这局面的逻辑迟早该把我引向把我找到的东西传达出去。既然直到去年,我根本不肯为我的数学之爱付出那种大规模的、长期的投入——那本是必要的:得靠一番因事制宜的“逐件加工”,才能把我掘出的矿藏“开采”出来拿去发表——那么剩下的另一条路,就是把我至少最放在心上的那些东西,传达给几位足够“在行”的数学家朋友。

我想,要是在最近十年里我碰上一位数学家朋友,对我扮演对话者和消息来源的角色(就像五六十年代那许多年里,塞尔在很大程度所扮演的那样),同时又充当一道中转,把我能传给他的“消息”传出去(这角色从前塞尔不必扮演,因为我自己就担了!)——那么我想“在数学上有所行动”的欲望,就会得到足够的满足,足以化解我的落空,而我只消在数学上断断续续、有节制地投些精力,把最大的份额留给我新的爱。我第一次带着这样的期待(至少心里暗含这期待)去找一位数学家朋友,是1975年,最后一次是1982年,一年半以前。有趣的巧合:两次都是为了想办法“安置”(好让它传开去,谁知道呢,说不定最终还会有人发展它!)同一个“纲领”,同调与同伦代数方面的,其最初的胚芽要上溯到五十年代,而早在六十年代末以前就已经完全“成熟”(依照我内心的确信);这纲领的一个初步的、粗线条的展开,恰恰就是《追寻栈》的主题——而我此刻按说正该给它写导论!不管怎样,由于大概各不相同的缘由,我那些想重新找回一种“特权对话者”关系的尝试,都夭折了——那种关系从前(1970年以前)曾与塞尔有过,后来又与德利涅有过。不过有一个共同的情形:我准备给数学的时间相对有限。在我说的那两次(1975年和1982年),这肯定也促成了交流的瘸腿。其实,我主要想的是把什么东西“安置”出去,并不太费心去做那件必要的事——“(重新)跟上行情”,以便我这方面成为一个让通信者满意的对话者;而对方在行内那些通行技术上比我“在行”得多(至少可以这么说!)。

% ---------------- p.228 ----------------
\origpage{136} 我大可以把《致……的信》——它充当《追寻栈》的第一章,去年二月写的信,才一年多以前——看作我最后一次尝试:在我一位旧日朋友那里,为我如今某些想法与关切找到一点回响。从这封信开始(更确切地说,重新开始)的反思,后来竟要成为1970年以来第一篇注定要发表的数学文本——而这一点,我当时好几个星期都没有察觉。差不多一年之后,我才收到对这封内容充实的信的间接反应(参见注\fnoteP{2}{(38)}{这些笔记其实是那封写给……的长信的续篇,后来那封信成了第一章。它们打字打出来,是为了让这位旧日的朋友读得下去,另外还为了两三个人(尤其是龙尼·布朗[Ronnie Brown]),我以为他们会有兴趣。这封信其实从未得到回音,收信人也没读过它:将近一年之后(我问起他是否收到了),他真心感到惊讶——我居然会有一刻以为他可能读它,看看人们该指望从我这儿得到的是什么样的数学……} (38))。这反应比迄今任何一位数学同行写给我的信都更能说明问题,让我感到某些针对我这点微末之身的姿态——自我离开那个曾与他们共处的圈子以来,这些姿态在我的数学家朋友中间已经成了常事。这封信出自一位我曾当作朋友、怀着热忱的同情去信的人,而信里有一种蓄意的嘲弄,格外猛烈地让我想起一件事——那是我最近几年终于越来越清楚地意识到的事。此前,我主要有机会注意到的是对我这个人本身的疏远,在数学“大世界”里,首先是在那些曾与我或远或近算朋友的人中间(45)。那里说的不再是人际层面上的疏远,而是一种共识,其性质如同一种时尚,并且像时尚那样把自己摆成理所当然的东西,流行于一些多少“在行”的人之间:那种论千页一包一包的数学,还有那些我拿它们在一二十年里磨破了人耳朵的概念(46,47),全面权衡起来并不很严肃;那里头有大量装腔作势的轰响,值得的东西却没多少;而除开围绕概形与平展上同调这两个概念的那几大篇“普遍胡话”——它们有时候倒确实自有其用处;唉,这一点人还是乐意承认的——至少剩下的部分,忘掉更厚道;至于那些不顾雅趣、不顾显而易见的严肃准则,仍旧摆出要吹这种格罗滕迪克式小号的人,就该跟他们那位师父(不管认不认)扔进同一条口袋;他们要是受到应得的对待,也只能怪自己……

\origpage{137} 无疑,自1976年以来(50)、尤其是最近两三年来传到我这里的许多这一路的声音(我刚把它们“照原样”抄了下来),终于唤醒了我身上那根好斗的筋——最近十年里它多少睡过去了。它们像反射一样挑起了一个念头:投入混战,让那些什么也没弄懂的毛头小子闭上嘴——总之是个十足愚蠢的反射,就像公牛:只要给它看一块红布,在它鼻子前晃一晃,它立刻就抖擞起来、冲上去,忘了自己原本正安安静静走着的路,那可是它自己的路!不过我仍觉得这反射相当表皮,单靠它不足以让我动摇。再说,幸好,做数学显然比朝一块红布冲过去、四面挨刀要有滋味得多。可是做数学,而且不顾一切地坚持一种工作风格、一种看事情的路数——那些是我的——这也有点“投入混战”的意思;这是在那些轻蔑与拒斥的迹象面前挺立自己——那些迹象落到我身上,毫无疑问是回应我的旧友们在我身上感到、或者以为感到的轻蔑,那轻蔑若不是针对他们,至少也是针对他们毫无保留地认同的那个圈子。所以这多多少少也是在跟着那块红布跑,而不是走自己的路。

% ---------------- p.229 ----------------
\origpage{138} 这个念头,最近几周里已多次在我这里出现过;而今天这场反思,也许主要就是朝考察这一面去的。走着走着,又显出另一面,那里自我的力量无疑也占了大份,但它不像是一种单纯的好斗反射。不如说是一种我身上的欲望——此刻我还看不清它的本性——要给我最近十年或十二年所做的数学工作一个意义,或者要看到它取得它全部的意义;而这意义(我内心确信)不可能归结为私下里的乐趣或个人的历险。可是,即便这欲望的本性仍未被理解,既然我还没闲下来细看它,这场反思已足以向我表明:真正压在我身上、可以说硬掰着我的手让我往数学投入的那股力量,恰恰就在这欲望里——“翻转”之力。红布有没有,它照样起作用。如果它是依恋一段过去的迹象,那也是最近这十年的过去,也就是“1970年以后”的过去,而不是那些已经白纸黑字写下的东西、已经做成的东西的过去,不是1970年以前的。

\origpage{139} 骨子里,对这些东西、对将来“后世”会怎样处置它们(何况是否真有后世还很可疑……),我心里毫无不安。我在这段过去里感兴趣的,绝不是我在其中做过的事(以及它已有或将会有的运道),而恰恰是没有做成的事——在那份当年摆在我眼前的大纲领里,只有极小一部分,靠我的努力,以及靠那些有时愿意加入我的友人和学生的努力,得以实现。这份纲领本身,在我对数学事物的眼光与路数更新的同时,也未经预谋、未经寻求地更新了。年复一年,重心移动了,题目如此,我要说的东西本身也如此:如今我最首要的意图,不再是完成那些精细基础的大工程,而是去探那些最让我着迷的奥秘,比如“Motif”的奥秘,或者 $\mathbb{Q}$ 在 $\mathbb{Q}$ 上的伽罗瓦群的“几何”描述的奥秘。当然,一路走来,我至少免不了在这里那里勾画些基础,就像我已经开始做的那样(除了别的)在《穿越伽罗瓦理论的长征》里,或者像我现在正在《追寻栈》里做的那样。然而要说的东西变了,表达它的风格也变了。

% ---------------- p.230 ----------------
换个说法:最近十年里,在数学事物的世界里,我瞥见了一些神秘的、极美的东西。这些东西不是我私人的,它们生来就是要被传达的——瞥见它们这件事本身的意义,我这样感到,就在于把它们传达出去,好让人接过去、理解、吸收……然而传达它们,哪怕是向自己传达,也就是深入它们,多少发展它们——那是一件工作。我当然很清楚:我绝不会把这工作做到头,即便再给我一百年去投入。但这不该是我今天操心的事——在我余下活着、余下发现世界的时间里,我要拿多少年、多少个月投到那件工作上;而另一件工作正等着我,那是我唯一能做的。我生命的季节,不在我的权柄之内,也不是我的角色去安排的。
